\begin{table}[ht]
\centering
\small
\caption{Justification of the normal quantile in Equation~(\ref{eq:overshoot-threshold}), restricted to the two-sided regularity region $m'\min(p_0,1-p_0)\geq 10$. Because $\hat\phi_N$ is strictly decreasing in the count $K$, the rule $\hat\phi_N < \phi_1$ is \emph{identical} to the cut $K > m'\cos^2(N\phi_1)$, so the only quantity that has to be Gaussian is $K$ itself, and $K$ is a sum of $m'$ i.i.d.\ Bernoulli variables. The Berry--Esseen theorem then bounds the error of its normal approximation by $C(p_0^2+(1-p_0)^2)/\sqrt{m'p_0(1-p_0)}$ with $C \leq 0.4748$, non-asymptotically and for every $m'$; the second column reports the worst case of that bound and the third the distance actually attained. The remaining columns give the largest gap between the size the rule achieves and its nominal level $\alpha = 1-\mathrm{conf}$, computed from the exact binomial over safe values of $N$. That gap includes finite-sample threshold placement and binomial discreteness; its worst case need not shrink here because the regularity region expands toward the boundary as $m'$ grows.}
\label{tab:overshoot-operating}
\setlength{\tabcolsep}{6pt}
\begin{tabular}{r cc ccc}
\toprule
$m'$ & \makecell{Berry--Esseen \\ bound} & \makecell{Exact Kolmogorov distance \\ to normal approximation} & \makecell{$\mathrm{conf} = 0.5$ \\ ($\alpha = 0.5$)} & \makecell{$\mathrm{conf} = 0.66$ \\ ($\alpha = 0.34$)} & \makecell{$\mathrm{conf} = 0.95$ \\ ($\alpha = 0.05$)} \\
\cmidrule(lr){2-3}\cmidrule(lr){4-6}
& \multicolumn{2}{c}{normal approximation to $K$} & \multicolumn{3}{c}{largest deviation of the achieved size} \\
\midrule
50 & 0.114 & 0.084 & 0.077 & 0.077 & 0.037 \\
200 & 0.139 & 0.083 & 0.081 & 0.067 & 0.036 \\
800 & 0.147 & 0.083 & 0.080 & 0.072 & 0.040 \\
\bottomrule
\end{tabular}
\end{table}
